\section*{Chapter \ref{ch:ll:cpp-for-latency-ii-compile-time} --- C++ for Latency II: Compile Time}

\begin{solution}{exo:ll:cpp-for-latency-ii-compile-time:1}
On the committed measurement, $1.94/0.52 \approx 3.7$ times on one message type and $8.69/5.22 \approx 1.7$ times on the
mix.
\end{solution}

\begin{solution}{exo:ll:cpp-for-latency-ii-compile-time:2}
$(7.26 - 1.14)\ \text{ns} \times 5 \times 10^6 \approx 31$ milliseconds of processing a second: 3\% of a core, and six
nanoseconds on every event's latency.
\end{solution}

\begin{solution}{exo:ll:cpp-for-latency-ii-compile-time:3}
$250 \times 10^2 = 25\,000$; $2\,500 / 10^2 = 25$. \texttt{rescale(1, 2, 20)}: the target scale is out of range, so the
unlikely branch calls the cold reporting function and returns 0.
\end{solution}

\begin{solution}{exo:ll:cpp-for-latency-ii-compile-time:4}
The \texttt{switch} is compiled to a jump table (an indirect jump) or a chain of conditional branches; either depends on the
message type, and on a random mix the predictor is wrong about three times in four.
\end{solution}

\begin{solution}{exo:ll:cpp-for-latency-ii-compile-time:5}
Stage 1 (add the quantity) is the load of \texttt{qty} and the \texttt{add}; stage 2 the \texttt{xor} with the price and the
store of \texttt{acc}; stage 3 the \texttt{cmp} with the limit and the \texttt{jge}; stage 4 the \texttt{add} to the sum in
memory. The objects of stages 1 and 2 have disappeared: they have no data members, only code, which was inlined; stage 3's
limit and stage 4's sum remain as loads and stores at offsets of the pipeline object.
\end{solution}

\begin{solution}{exo:ll:cpp-for-latency-ii-compile-time:6}
Make the engine a template over the strategy type and instantiate it for each of the three; at start-up, read the
configuration and call the \texttt{run} of the matching instantiation once. The single dispatch happens outside the hot
loop.
\end{solution}

\begin{solution}{exo:ll:cpp-for-latency-ii-compile-time:7}
The listing gains a \texttt{call} to the stage, with the event's fields stored to memory before it and reloaded after, since
the compiler can no longer see what the callee changes; the time per event rises toward the dynamic pipeline's, without the
indirect load of a table.
\end{solution}

\begin{solution}{exo:ll:cpp-for-latency-ii-compile-time:8}
The benchmark measures a perfectly predicted call in a tight loop. It misses the mispredictions when targets vary, and above
all the lost optimisation: a virtual call cannot be inlined, so values are spilled around it and nothing is folded across
it. The four-stage pipeline lost six nanoseconds, not four calls' worth of 1.9.
\end{solution}

\begin{solution}{pb:ll:cpp-for-latency-ii-compile-time:1}
\begin{enumerate}
\item One type: virtual 1.9, variant 1.8, CRTP \qty{0.5}{\nano\second}; mix: 8.7, 6.4 and \qty{5.2}{\nano\second} (committed
measurement).
\item \omterm{def:ll:cpu-microarchitecture:branch}{Branch misprediction} on the message type, which chooses the target.
\item About 1.1 and \qty{7.3}{\nano\second} per event.
\item Its stages run in the same order for every event; the only data-dependent branch is the gate, the same in both.
\item It loads the next stage pointer, loads the object's table, calls through it, tests the result and loops; registers
are saved around each call.
\item Thirteen instructions.
\item Folding the stages' arithmetic together, keeping the event in registers, removing the empty stage objects, and a
single conditional branch for the gate.
\item It lets the compiler devirtualise calls through references of the concrete type; it does nothing for calls through
the base class pointer that a plug-in architecture uses.
\item Make each plug-in a template parameter and instantiate the engine for each supported combination; choose the
instantiation once at start-up.
\item New plug-ins cannot be loaded into a running binary; adding one needs a rebuild and a release.
\item $3 \times 2 = 6$.
\item Each variant is built and runs the full test and replay suite; a matrix in continuous integration.
\item \textbf{Named result.} Virtual dispatch costs about 3.7 times CRTP when the message type is predictable and about 1.7
times on a random mix; the four-stage pipeline costs about \qty{1.1}{\nano\second} per event composed statically and
\qty{7.3}{\nano\second} through virtual calls.
\item It measures the call alone, predicted, and not what the call prevents: \omterm{def:ll:cpp-for-latency-ii-compile-time:inline}{inlining} and the optimisations it enables.
\item Off the hot path (configuration, reporting, tools), and where the set of types is open (plug-ins loaded without a
rebuild in a research engine).
\item A named requirement checked at the point of use: the error for a stage without \texttt{bool on(E\&)} is one line.
\item Tables and constants (scales, lengths, lookup tables) built without start-up code and checked by
\texttt{static\_assert}.
\item To keep the hot function small and contiguous in the instruction cache, with the common case falling through.
\item Because \texttt{always\_inline} is a request to a heuristic that can change between compiler versions; the assembly
shows what happened.
\item Everything that is fixed for the life of the process: types, strategies, tables.
\end{enumerate}
\end{solution}

\begin{solution}{iq:ll:cpp-for-latency-ii-compile-time:1}
The object holds a pointer to its class's table of function pointers; the call loads the pointer, loads the entry and calls
through it. Cost: two dependent loads and an indirect branch (cheap when predicted), and the loss of \omterm{def:ll:cpp-for-latency-ii-compile-time:inline}{inlining}.
\iqlookfor{the vtable mechanism and the indirect cost.}
\end{solution}

\begin{solution}{iq:ll:cpp-for-latency-ii-compile-time:2}
A base class template parameterised by its derived class calls the derived class's functions through
\texttt{static\_cast<Derived\&>(*this)}; dispatch is resolved at compile time and inlinable. Use it for a fixed family of
behaviours on a hot path (handlers, policies).
\iqlookfor{compile-time dispatch and its trade-off: no runtime substitution.}
\end{solution}

\begin{solution}{iq:ll:cpp-for-latency-ii-compile-time:3}
A \texttt{constexpr} function may be evaluated at compile time when its arguments are constants, and otherwise runs at run
time; a \texttt{consteval} function must be evaluated at compile time, so a call that cannot be is an error.
\iqlookfor{may versus must.}
\end{solution}

\begin{solution}{iq:ll:cpp-for-latency-ii-compile-time:4}
It exposes the callee's body to the caller's optimisations: constants propagate, values stay in registers, dead work
disappears. It can make code slower by growing it beyond the instruction cache or the loop buffer, or by \omterm{def:ll:cpp-for-latency-ii-compile-time:inline}{inlining} rare code
into a hot loop.
\iqlookfor{optimisation across the boundary, and code size.}
\end{solution}

\begin{solution}{iq:ll:cpp-for-latency-ii-compile-time:5}
A variant: the set of message types is closed, the object needs no heap allocation, dispatch is a jump on the index, and
the handlers can be inlined. A hierarchy only if the set must be open.
\iqlookfor{closed versus open sets of types, and allocation.}
\end{solution}

\begin{solution}{iq:ll:cpp-for-latency-ii-compile-time:6}
The engine is a template over its strategy (and other policies); each supported strategy is an instantiation compiled into
the binary; the configuration selects one instantiation at start-up with a single dispatch outside the loop; the hot loop
is fully static. Tested as a matrix.
\iqlookfor{one dispatch at start-up, static thereafter.}
\end{solution}
